\documentclass[10pt,a4paper]{amsart}
\usepackage[margin=19mm]{geometry}
\usepackage{amsmath,amssymb,mathtools}
\usepackage[scr=rsfs,cal=euler]{mathalfa}
\usepackage{xfrac}
\usepackage[unicode,hidelinks,bookmarks=false]{hyperref}
\newcommand{\ie}{i.e.\ }
\newcommand{\cf}{cf.\ }
\newcommand{\resp}{respectively\ }
\newcommand{\Subsection}{\subsection}
\providecommand{\nobreakdash}{\nobreak\hbox{-}}
\setlength{\emergencystretch}{2em}
\multlinegap0pt
\usepackage{amssymb}
\usepackage{mathtools}
\usepackage[shortlabels]{enumitem}
\usepackage{xcolor}
\newtheorem{lemma}{Lemma}
\newtheorem{corollary}{Corollary}
\newtheorem{proposition}{Proposition}
\newcommand{\C}{\mathbb C}
\newcommand{\MM}{\mathcal M}
\newcommand{\N}{\mathbb N}
\newcommand{\PP}{\mathbb P}
\newcommand{\Z}{\mathbb Z}
\newcommand{\ord}{\operatorname{ord}}
\title{Meromorphic solutions of first-order differential equations with rational exponential coefficients}
\author{Zhong Deguang, Meng Fanning\textsuperscript{*}, Yuan Wenjun}
\date{Source: arXiv:2609.10219v1 (CC BY 4.0)}
\begin{document}
\maketitle

\noindent\emph{Source and licence.} Meromorphic solutions of first-order differential equations with rational exponential coefficients. Version arXiv:2609.10219v1 (CC BY 4.0), distributed by arXiv under the Creative Commons Attribution 4.0 International licence, http://creativecommons.org/licenses/by/4.0/, as recorded in the arXiv licence metadata for this exact version and shown on the abstract page https://arxiv.org/abs/2609.10219v1. Full licence notice: \texttt{licenses/arxiv-2609.10219-CC-BY-4.0.txt}.\par\smallskip

\noindent\textit{Editorial scope.} Counted unit: the article's proposition prop:comparison (source0.tex lines 798--807). The article's complete original proof of it is delivered with it (source0.tex lines 809--917, one \texttt{proof} environment of 109 lines). No mathematics is written, repaired or re-derived: the statement and the proof are the article's own lines, in the article's own notation, and no retained line is substituted or re-flowed.  Retained uncounted as the source-local context the proof needs: lines 278--286; lines 386--395; lines 440--444; lines 590--596; lines 729--742.  Nothing was omitted from any retained span, so the package carries no omission record.  No retained line contains a citation, so the wrapper carries no bibliography and no bibliography entry.  No retained line contains a figure environment, an image-inclusion command or drawing code, so no image is delivered and the package declares no asset dependency.  Editorial formatting only, in addition: an AMS \texttt{amsart} preamble that restates the article's own declarations and macros used by the retained lines, namely the environments \texttt{lemma}, \texttt{corollary}, \texttt{proposition} on separate counters, together with the standard packages those lines need.  The retained environments print in this document as Lemma 1, Corollary 1, Proposition 1 in order of appearance, whereas the article prints the same items with the numbering of its own full document.  The complete native source of the article as distributed in the arXiv e-print for this exact version is bundled unchanged as \texttt{sources/209885/source0.tex}.
\begin{lemma}
\label{lem:relative}
If $x$ is transcendental over a field $F$, then $F$ is relatively
algebraically closed in $F(x)$. Consequently, if $v\in\C(t,w)$ and
\[
 v^d=\gamma t,\qquad \gamma\in\C^*,\quad d\in\N,
\]
then $d=1$.
\end{lemma}

\begin{lemma}[Regular derivations on a rational curve]
\label{lem:riccati}
Let $K$ be a differential field and let $x$ be transcendental over $K$.
For an extension of the derivation to $K(x)$, the following conditions
are equivalent:
\begin{enumerate}
\item $D$ is regular on $\PP^1_{\overline K}$;
\item $Dx=a_2x^2+a_1x+a_0$, where $a_0,a_1,a_2\in K$.
\end{enumerate}
\end{lemma}

\begin{corollary}
\label{cor:change-generator}
Suppose that $K(x)=K(y)$ and $Dy$ is a polynomial of degree at most two
in $y$. Then $Dx$ is a polynomial of degree at most two in $x$.
\end{corollary}

\begin{lemma}
\label{lem:autonomous}
Let $H\in\MM(\C)$ be nonconstant and satisfy $H'=b(H)$ with
$b\in\C(s)$. A constant-coefficient fractional linear change of the
dependent variable transforms $H$ into either $z$ or $e^{\lambda z}$,
where $\lambda\in\C^*$.
\end{lemma}

\begin{proposition}
\label{prop:two-point}
Let $K=\C(h)$ be relatively algebraically closed in a field $L$ of
transcendence degree one over $K$. Suppose that $L$ is the function
field of a smooth projective curve $\Gamma/K$, that $D$ is regular on
its geometric curve, and that $L$ admits an injective differential map
to $\MM(\C)$. If $h,t$ are algebraically independent,
$Dt=t$ and $\iota(t)=e^z$, then there exist
$u\in L$, $d\in\N$ and $c\in K^*$ such that
\begin{equation}
 L=K(u),\qquad t=c(h)u^d.
 \label{eq:kummer-form}
\end{equation}
\end{proposition}

\begin{proposition}
\label{prop:comparison}
Let $L=\C(t,w)$ with $t,w$ algebraically independent. Let $D$ be a
derivation with $Dt=t$, and suppose that an injective differential
homomorphism $\iota:L\to\MM(\C)$ satisfies $\iota(t)=e^z$.
Assume that there is a subfield $K=\C(h)$ relatively algebraically
closed in $L$ such that $D(K)\subset K$ and $D$ is regular on the
geometric generic curve of $L/K$. Then $Dw$ is a polynomial of degree
at most two in $w$, with coefficients in $\C(t)$.
\end{proposition}

\begin{proof}
We distinguish whether the two base functions are algebraically
dependent.

\smallskip\noindent
\emph{The dependent case.}
If $h$ is algebraic over $\C(t)$, Lemma~\ref{lem:relative} implies
$h\in\C(t)$. Since $h$ is nonconstant, $t$ is then algebraic over
$\C(h)$. Relative algebraic closedness of $K$ gives $t\in K$, hence
\begin{equation}
 K=\C(t).
 \label{eq:equal-bases}
\end{equation}
The smooth projective curve with field $L/K=\C(t)(w)$ is
$\PP^1_K$. Regularity and Lemma~\ref{lem:riccati} give the desired
conclusion in the original coordinate $w$.

\smallskip\noindent
\emph{The independent case.}
Suppose that $h,t$ are algebraically independent. By
Proposition~\ref{prop:two-point},
\begin{equation}
 L=\C(h,u),\qquad t=c(h)u^d.
 \label{eq:comparison-kummer}
\end{equation}
Put $H=\iota(h)$ and $U=\iota(u)$. Both are global meromorphic
functions. Lemma~\ref{lem:autonomous} permits a change of the base
coordinate so that $H=z$ or $H=e^{\lambda z}$. Such a change leaves
$K$ unchanged and merely rewrites the rational function $c$.

If $H=z$, substitution in \eqref{eq:comparison-kummer} gives
\begin{equation}
 e^z=c(z)U(z)^d.
 \label{eq:additive-divisor-identity}
\end{equation}
At every $a\in\C$, the left side has order zero. Therefore
\[
 \ord_a c+d\,\ord_a U=0.
\]
In particular, every finite zero and pole of $c$ has order divisible
by $d$. Factoring this rational function into linear factors yields
\begin{equation}
 c(h)=c_0r(h)^d,
 \qquad c_0\in\C^*,\quad r\in\C(h)^*.
 \label{eq:additive-divisibility}
\end{equation}
Set $v=r(h)u$. Then $L=\C(h,v)$ and $t=c_0v^d$.
Lemma~\ref{lem:relative} gives $d=1$. Consequently
\begin{equation}
 L=\C(t,h),\qquad Dh=1.
 \label{eq:additive-product}
\end{equation}
Use Corollary~\ref{cor:change-generator} over $\C(t)$, with $h$ as
the rational generator. It follows that $Dw$ is of Riccati type.

Suppose instead that $H=e^{\lambda z}$, $\lambda\ne0$. For every
$a\in\C^*$ there is a time $z_0$ with $H(z_0)=a$, and
$H'(z_0)=\lambda a\ne0$. Comparing orders in
\[
 e^z=c(H(z))U(z)^d
\]
at this time proves
\begin{equation}
 d\mid\ord_a c\qquad(a\in\C^*).
 \label{eq:multiplicative-divisibility}
\end{equation}
The two omitted base values are precisely the ones at which this
argument imposes no separate divisibility condition. Factorization in
$\C(h)$ now gives
\begin{equation}
 c(h)=c_0h^k r(h)^d,
 \qquad c_0\in\C^*,\quad k\in\Z,
 \quad r\in\C(h)^*.
 \label{eq:monomial-coefficient}
\end{equation}
With $v=r(h)u$ we obtain
\begin{equation}
 L=\C(h,v),\qquad t=c_0h^kv^d.
 \label{eq:monomial-map}
\end{equation}
Injectivity gives $Dh=\lambda h$. Differentiating the monomial
identity gives
\begin{equation}
 Dv=\mu v,\qquad \mu=\frac{1-k\lambda}{d}.
 \label{eq:diagonal-derivation}
\end{equation}

Let $g=\gcd(k,d)$, taken positive. If $g>1$, then the element
$h^{k/g}v^{d/g}\in L$ has $g$th power $t/c_0$. This contradicts
Lemma~\ref{lem:relative}. Thus $\gcd(k,d)=1$.
Choose integers $a,b$ with
\begin{equation}
 kb-da=1,
 \label{eq:bezout}
\end{equation}
and put $y=h^av^b$. The inverse change of variables is explicit:
\begin{equation}
 h=(t/c_0)^b y^{-d},\qquad
 v=(t/c_0)^{-a}y^k.
 \label{eq:inverse-monomial}
\end{equation}
Hence
\begin{equation}
 L=\C(t,y),\qquad Dy=(a\lambda+b\mu)y.
 \label{eq:original-base-product}
\end{equation}
Corollary~\ref{cor:change-generator}, again over the original base
$\C(t)$, proves the assertion for $Dw$.
\end{proof}

\end{document}
